\documentclass[12pt,letterpaper]{article}
\usepackage[utf8]{inputenc}
\usepackage[T1]{fontenc}
\usepackage[spanish,es-nodecimaldot]{babel}
\usepackage{newtxtext,newtxmath}
\usepackage[margin=1in,headheight=15pt]{geometry}
\usepackage{setspace,graphicx,booktabs,tabularx,array,float,caption,fancyhdr,xurl,hyperref,titlesec,listings,amsmath}
\hypersetup{hidelinks}\urlstyle{same}
\doublespacing\setlength{\parindent}{0.5in}\setlength{\parskip}{0pt}
\setcounter{secnumdepth}{0}
\titleformat{\section}{\normalfont\bfseries\centering}{}{0pt}{}
\titleformat{\subsection}{\normalfont\bfseries}{}{0pt}{}
\titlespacing*{\section}{0pt}{1em}{0pt}
\addto\captionsspanish{\renewcommand{\tablename}{Tabla}}
\captionsetup{justification=raggedright,singlelinecheck=false,labelsep=newline,labelfont=bf,textfont=it}
\pagestyle{fancy}\fancyhf{}\rhead{\thepage}\renewcommand{\headrulewidth}{0pt}
\fancypagestyle{plain}{\fancyhf{}\rhead{\thepage}\renewcommand{\headrulewidth}{0pt}}
\newcommand{\nota}[1]{\par\noindent\small\textit{Nota.} #1\par\normalsize}
\newcommand{\refapa}[1]{\par\noindent\hangindent=0.5in\hangafter=1 #1\par}
\lstset{language=Python,basicstyle=\ttfamily\scriptsize,breaklines=true,showstringspaces=false,frame=single,numbers=left,numberstyle=\tiny}
\begin{document}\raggedright
\begin{titlepage}\thispagestyle{fancy}\begin{center}
\vspace*{1in}{\bfseries Reconocimiento de dígitos manuscritos con redes neuronales densas: análisis comparativo de algoritmos de optimización, regularización y generalización\par}\vspace{1\baselineskip}
Karen Sofía Delgado Calderon\par
Especialización en Inteligencia Artificial\par Universidad de Cundinamarca\par
608142 -- Deep Learning - Conceptos\par R1-A3-S7 -- Redes Neuronales III\par 10 de octubre de 2026\par
\end{center}\vfill\end{titlepage}\setcounter{page}{2}
\begin{center}\textbf{Reconocimiento de dígitos manuscritos con redes neuronales densas: análisis comparativo de algoritmos de optimización, regularización y generalización}\end{center}

\section{Planteamiento del problema y fundamentación experimental}
El reconocimiento automático de caracteres manuscritos representa un problema canónico en la visión por computador y el aprendizaje estadístico de patrones (Alpaydin \& Kaynak, 1998; LeCun et al., 1998). Aunque el estado del arte contemporáneo privilegia arquitecturas convolucionales y mecanismos de atención para procesar imágenes continuas, la investigación de modelos basados en perceptrones multicapa densos (\textit{Multilayer Perceptron}, MLP) sobre cuadrículas de píxeles vectorizadas constituye un escenario idóneo para analizar el impacto intrínseco de los algoritmos de optimización y la regularización estocástica sin la interferencia del sesgo inductivo espacial de los filtros convolucionales (Bishop, 2006; Goodfellow et al., 2016).

En esta investigación se utilizó el conjunto de datos de referencia \textit{Optical Recognition of Handwritten Digits} (\textit{Digits}; Alpaydin \& Kaynak, 1998; Pedregosa et al., 2011; scikit-learn, s.~f.). La colección está conformada por 1797 mapas de bits normalizados de $8 \times 8$ píxeles en escala de grises correspondientes a los diez dígitos arábigos (clases del 0 al 9). Cada instancia está descrita por un vector unidimensional de 64 atributos continuos con intensidades comprendidas en el rango entero $[0, 16]$, derivadas del conteo de píxeles encendidos en bloques de muestreo originales de $4 \times 4$ de una matriz primaria de $32 \times 32$.

Para garantizar una evaluación rigurosa y prevenir cualquier fenómeno de fuga de información (\textit{data leakage}), la totalidad de las 1797 instancias se segmentó mediante muestreo estratificado en tres particiones disjuntas con proporciones $60\,\%$, $20\,\%$ y $20\,\%$, correspondientes a 1077 observaciones de entrenamiento, 360 de desarrollo (\textit{dev/validation}) y 360 de prueba independiente (\textit{test}). El preprocesamiento consistió en un reescalado lineal invariante:
\begin{equation}
x_{\text{norm}} = \frac{x}{16{,}0}
\end{equation}
A diferencia de transformaciones que estiman momentos muestrales de media y desviación típica sobre los datos, la división por $16{,}0$ responde a una cota física intrínseca del hardware del escáner y la compresión por bloques definida por Alpaydin y Kaynak (1998). Por consiguiente, este escalado es determinista, libre de parámetros aprendidos y garantiza inmunidad absoluta contra fugas informativas hacia las particiones de desarrollo y prueba.

La arquitectura de red neuronal configurada comprende una secuencia completamente densa $64 \to 128 \to 64 \to 10$, totalizando 17226 parámetros entrenables ($[64\times 128 + 128] + [128\times 64 + 64] + [64\times 10 + 10]$). Las capas intermedias incorporan funciones de activación ReLU (Nair \& Hinton, 2010) e inicialización uniforme de Glorot y Bengio (2010). La capa de salida emite 10 \textit{logits} sobre los cuales se evalúa analíticamente la función de pérdida de entropía cruzada multiclase (Paszke et al., 2019):
\begin{equation}
\mathcal{L}_{\text{CE}}(y, z) = -\sum_{c=0}^{9} y_c \log \left( \frac{\exp(z_c)}{\sum_{j=0}^{9} \exp(z_j)} \right)
\end{equation}
La Tabla~\ref{tab:protocolo} compendia las especificaciones metodológicas del protocolo experimental adoptado.

\begin{table}[H]\caption{Protocolo experimental y partición de datos para el reconocimiento de dígitos}\label{tab:protocolo}
\begin{singlespace}\small\centering\begin{tabularx}{\textwidth}{>{\raggedright\arraybackslash}p{0.36\textwidth} X}\toprule
Elemento de diseño & Especificación metodológica \\\midrule
Conjunto de datos & Optical Recognition of Handwritten Digits: 1797 muestras, 64 píxeles, 10 dígitos \\
Particiones de datos & Entrenamiento: 1077; Desarrollo (\textit{dev}): 360; Prueba (\textit{test}): 360 instancias \\
Estrategia de división & Muestreo estratificado multiclase balanceado; prueba estrictamente ciega \\
Preprocesamiento & Escalamiento lineal determinista invariante ($x / 16{,}0$), sin fuga muestral \\
Semillas pseudoaleatorias & 11, 23, 37 (reproducibilidad estricta de pesos en CPU y tensores) \\
Régimen de entrenamiento & Minilotes de 64 muestras, tope máximo de 100 épocas \\
Criterio rector de selección & Mayor $F_1$ macro promedio en el conjunto de desarrollo \\\bottomrule
\end{tabularx}\end{singlespace}\nota{Diseño experimental concebido para la práctica R1-A3-S7 de la Especialización en Inteligencia Artificial. Los datos provienen del repositorio de scikit-learn (s.~f.), basados en el estándar de Alpaydin y Kaynak (1998).}
\end{table}

\section{Dinámica de optimización y mecanismos de regularización}
El entrenamiento de redes profundas implica resolver un problema de optimización no convexo donde coexisten mínimos locales, curvaturas anisotrópicas y puntos de ensilladura (Goodfellow et al., 2016; Ruder, 2016). En este experimento se evaluaron de forma sistemática tres familias de algoritmos de optimización: Descenso de Gradiente Estocástico con momento (\texttt{SGD\_momentum}; Polyak, 1964; Bottou, 2010), RMSprop con media móvil cuadrática de gradientes (Tieleman \& Hinton, 2012) y Adam con corrección de sesgo para momentos de primer y segundo orden (Kingma \& Ba, 2015).

Bajo SGD con momento clásico de Polyak (1964), la variable de velocidad acumula los gradientes de minilotes anteriores con un factor de persistencia $\beta = 0{,}9$, amortiguando oscilaciones transversales y acelerando la navegación en direcciones con gradiente persistente:
\begin{equation}
v_t = \beta v_{t-1} + \eta \nabla_\theta \mathcal{L}(\theta_{t-1}), \quad \theta_t = \theta_{t-1} - v_t
\end{equation}
Por otro lado, RMSprop (Tieleman \& Hinton, 2012) normaliza las actualizaciones individuales dividiendo el gradiente por la raíz de la media móvil exponencial de gradientes al cuadrado, permitiendo que las coordenadas con gradientes tenues reciban pasos efectivos más amplios. Finalmente, Adam (Kingma \& Ba, 2015) sintetiza ambas ideas combinando estimaciones corregidas por sesgo del primer momento ($\hat{m}_t$) y del segundo momento ($\hat{v}_t$):
\begin{equation}
\theta_t = \theta_{t-1} - \frac{\eta}{\sqrt{\hat{v}_t} + \epsilon} \hat{m}_t
\end{equation}

Con el objeto de evaluar el impacto específico de la regularización en la capacidad de generalización frente a la sobreparametrización (17226 parámetros frente a 1077 ejemplos de entrenamiento; Bishop, 2006; Hastie et al., 2009), se contrastaron cinco variantes algorítmicas, descritas en la Tabla~\ref{tab:config}.

\begin{table}[H]\caption{Configuraciones experimentales y variantes de regularización ensayadas}\label{tab:config}
\begin{singlespace}\small\centering\begin{tabularx}{\textwidth}{lcccc}\toprule
Configuración algorítmica & Tasa base ($\eta$) & Decaimiento ($\lambda$) & \textit{Dropout} ($p$) & Parada temprana \\\midrule
\texttt{SGD\_momentum} & 0,0600 & 0,000 & 0,00 & Inactiva \\
\texttt{Adam\_base} & 0,0030 & 0,000 & 0,00 & Inactiva \\
\texttt{RMSprop} & 0,0015 & 0,000 & 0,00 & Inactiva \\
\texttt{Adam\_decay} & 0,0030 & 0,001 & 0,00 & Inactiva \\
\texttt{Adam\_completo} & 0,0030 & 0,001 & 0,20 & Activa (paciencia 16) \\\bottomrule
\end{tabularx}\end{singlespace}\nota{Total de 15 corridas experimentales (5 configuraciones bajo 3 semillas deterministas). Los parámetros de tasa se seleccionaron mediante ajuste previo de convergencia estable.}\end{table}

En la variante \texttt{Adam\_decay}, se añadió decaimiento de pesos clásico mediante penalización de norma $L_2$ con $\lambda = 0{,}001$ (Krogh \& Hertz, 1991). Por su parte, la variante \texttt{Adam\_completo} incorporó una estrategia de regularización compuesta que incluye \textit{dropout} estocástico con probabilidad $p = 0{,}20$ en las capas ocultas (Srivastava et al., 2014), un programador de tasa de aprendizaje adaptativo (\texttt{ReduceLROnPlateau}) que divide la tasa por 2 si la pérdida en desarrollo no mejora tras 6 épocas continuas, y parada temprana (\textit{early stopping}; Prechelt, 1998) con paciencia de 16 épocas. La mecánica estocástica de \textit{dropout} desactiva neuronas mediante una máscara de Bernoulli:
\begin{equation}
r^{(l)} \sim \text{Bernoulli}(1-p), \quad \tilde{y}^{(l)} = r^{(l)} \odot y^{(l)}
\end{equation}

El núcleo de optimización y propagación implementado en PyTorch (Paszke et al., 2019) se expone a continuación:

\begin{singlespace}\begin{lstlisting}
optimizador = torch.optim.Adam(modelo.parameters(), lr=0.003, weight_decay=0.001)
planificador = torch.optim.lr_scheduler.ReduceLROnPlateau(optimizador, mode='min', factor=0.5, patience=6)
for imagenes_batch, etiquetas_batch in cargador_entrenamiento:
    optimizador.zero_grad(set_to_none=True)
    salidas = modelo(imagenes_batch)
    perdida = nn.functional.cross_entropy(salidas, etiquetas_batch)
    perdida.backward()
    optimizador.step()
\end{lstlisting}\end{singlespace}

\section{Evidencia empírica en desarrollo y selección del modelo}
Las curvas de convergencia para la pérdida de entrenamiento y la macro-métrica $F_1$ en el conjunto de desarrollo se grafican en la Figura~\ref{fig:optimizadores} para la semilla 11. Se comprueba que los algoritmos adaptativos (Adam y RMSprop) exhiben una aceleración inicial superior a SGD con momento, logrando una reducción de la entropía cruzada por debajo de 0,20 antes de la época 30. La Tabla~\ref{tab:dev} compendia las métricas agregadas obtenidas tras promediar las tres ejecuciones deterministas independientes.

\begin{figure}[H]\caption{Evolución de pérdida y macro F1 en desarrollo según la configuración}\label{fig:optimizadores}
\centering\includegraphics[width=1.0\textwidth,height=0.45\textheight,keepaspectratio]{resultados/optimizadores.png}
\nota{Trayectorias de aprendizaje en la partición de desarrollo para la inicialización con semilla 11. El cálculo de $F_1$ macro pondera idénticamente los 10 dígitos sin sesgo de clase.}\end{figure}

\begin{table}[H]\caption{Rendimiento en desarrollo consolidado a través de las tres semillas}\label{tab:dev}
\begin{singlespace}\small\centering\begin{tabularx}{\textwidth}{lcccc}\toprule
Configuración algorítmica & $F_1$ macro medio & Desv. Est. (DE) & Pérdida media & Brecha $F_1$ (Train$-$Dev) \\\midrule
\texttt{SGD\_momentum} & 0,9641 & 0,0055 & 0,1816 & 0,0359 \\
\texttt{Adam\_base} & 0,9676 & 0,0016 & 0,1651 & 0,0324 \\
\texttt{RMSprop} & 0,9657 & 0,0065 & 0,1631 & 0,0321 \\
\texttt{Adam\_decay} & \textbf{0,9714} & \textbf{0,0016} & \textbf{0,0915} & \textbf{0,0286} \\
\texttt{Adam\_completo} & 0,9705 & 0,0016 & 0,0727 & 0,0280 \\\bottomrule
\end{tabularx}\end{singlespace}\nota{DE representa la desviación estándar muestral insesgada entre semillas 11, 23 y 37. La brecha $F_1$ cuantifica el diferencial entre entrenamiento y validación. Cifras derivadas de \texttt{experimentos\_dev.csv}.}\end{table}

Tal como se desprende de la Tabla~\ref{tab:dev}, la configuración que alcanzó el valor más alto en el criterio rector de selección fue \texttt{Adam\_decay}, con una puntuación media de $F_1 = 0{,}9714$, una dispersión intersemilla mínima ($\text{DE} = 0{,}0016$) y una brecha de generalización de apenas $0{,}0286$. 

Resulta de suma relevancia teórica reflexionar sobre el comportamiento de la variante \texttt{Adam\_completo}: a pesar de incorporar regularización múltiple y reducir la pérdida a $0{,}0727$, su $F_1$ macro medio se situó en $0{,}9705$, ligeramente inferior al de \texttt{Adam\_decay}. En imágenes de baja resolución ($8 \times 8$), la información morfológica que distingue trazos delgados (como la continuidad de un lazo cerrado en un 8 frente a un 1 o la curvatura de un 3 frente a un 9) descansa en combinaciones dispersas de apenas 2 o 3 píxeles adyacentes. El silenciamiento estocástico del 20\,\% de las neuronas mediante \textit{dropout} interrumpe estas vías latentes críticas, induciendo un leve subajuste (\textit{underfitting}) que perjudica la discriminación de casos limítrofes (Srivastava et al., 2014; Bishop, 2006). Por consiguiente, la contracción analítica de pesos vía penalización $L_2$ (\texttt{Adam\_decay}) demostró ser la estrategia regularizadora óptima.

\begin{figure}[H]\caption{Contracción de la brecha de generalización entre entrenamiento y desarrollo}\label{fig:regularizacion}
\centering\includegraphics[width=0.9\textwidth,height=0.47\textheight,keepaspectratio]{resultados/regularizacion.png}
\nota{Evolución comparada de pérdidas de entrenamiento frente a desarrollo entre la variante base y la política regularizada para la semilla 11, evidenciando la estabilización del error.}\end{figure}

La Figura~\ref{fig:regularizacion} ratifica que el decaimiento de pesos estabiliza la trayectoria de validación sin la oscilación inducida por el apagado aleatorio de unidades. En consecuencia, el modelo \texttt{Adam\_decay} fue formalmente seleccionado y congelado para la evaluación en prueba.

\section{Evaluación en prueba y análisis cualitativo de errores}
Una vez fijada la arquitectura y los pesos correspondientes a la configuración \texttt{Adam\_decay}, se desbloqueó el conjunto de prueba aislado de 360 imágenes. Se calcularon las métricas globales y multiclase de exactitud (\textit{accuracy}), precisión macro, exhaustividad (\textit{recall}) macro y $F_1$ macro ponderado (Sokolova \& Lapalme, 2009; Grandini et al., 2020). La Tabla~\ref{tab:test} detalla las evaluaciones para cada semilla determinista.

\begin{table}[H]\caption{Evaluación final en el conjunto de prueba independiente}\label{tab:test}
\begin{singlespace}\small\centering\begin{tabularx}{\textwidth}{cccccc}\toprule
Semilla & Exactitud (\textit{Accuracy}) & Precisión macro & Exhaustividad macro & $F_1$ macro & Pérdida de prueba \\\midrule
11 & 0,9667 & 0,9675 & 0,9663 & 0,9662 & 0,1234 \\
23 & 0,9667 & 0,9687 & 0,9663 & 0,9665 & 0,1405 \\
37 & 0,9667 & 0,9674 & 0,9662 & 0,9656 & 0,1123 \\\bottomrule
\end{tabularx}\end{singlespace}\nota{Evaluación realizada sobre 360 imágenes no vistas durante el entrenamiento ni la selección de hiperparámetros. Se aprecia una notable robustez y reproducibilidad métrica intersemilla.}\end{table}

La matriz de confusión presentada en la Figura~\ref{fig:matriz_test} ilustra la distribución predictiva del modelo seleccionado para la semilla 11. Con una exactitud del $96{,}67\,\%$, 348 de las 360 instancias fueron clasificadas con total acierto, situándose densamente en la diagonal principal.

\begin{figure}[H]\caption{Matriz de confusión obtenida sobre las 360 imágenes de prueba}\label{fig:matriz_test}
\centering\includegraphics[width=0.82\textwidth,height=0.45\textheight,keepaspectratio]{resultados/matriz_test.png}
\nota{Matriz de contingencia calculada con el modelo definitivo (\texttt{Adam\_decay}) en la semilla 11. Las filas representan la clase manuscrita real y las columnas la predicción de la red.}\end{figure}

Para comprender la etiología del 3,33\,\% de desaciertos (12 dígitos erróneos), la Figura~\ref{fig:errores_test} expone visualmente los patrones de activación de píxeles correspondientes a los casos fallidos. El análisis cualitativo revela que los errores no provienen de deficiencias del optimizador, sino de ambigüedades morfológicas extremas inherentes a la baja resolución de $8 \times 8$:
\begin{itemize}
\item Dígitos 8 cuyos trazos superiores carecen de intensidad luminosa suficiente, generando una firma vectorial indistinguible de un dígito 1 vertical.
\item Dígitos 3 cuyos bucles curvados colapsan espacialmente en la matriz comprimida, siendo interpretados por la red como dígitos 9 u 8.
\item Dígitos 5 con esquinas inferiores redondeadas clasificados erróneamente como 6.
\end{itemize}
Esta limitación subraya la necesidad de recurrir a capas convolucionales cuando se requiere capturar invariancia translacional y topológica local en problemas tipográficos de mayor resolución (LeCun et al., 1998).

\begin{figure}[H]\caption{Inspección visual de dígitos de prueba erróneamente clasificados}\label{fig:errores_test}
\centering\includegraphics[width=0.9\textwidth,height=0.42\textheight,keepaspectratio]{resultados/errores_test.png}
\nota{Visualización de las instancias de prueba con discrepancia predictiva en resolución nativa de $8\times8$, evidenciando colapso de trazos y ambigüedad manuscrita.}\end{figure}

\section{Conclusiones y reproducibilidad}
El estudio comparativo permitió confirmar que el optimizador Adam con penalización por decaimiento de pesos ($L_2$) alcanza un equilibrio sobresaliente en el compromiso sesgo-varianza para la clasificación de dígitos vectorizados. La evidencia demostró que la adición de \textit{dropout} no siempre es beneficiosa: en representaciones densas comprimidas, el ruido estocástico puede destruir activaciones morfológicas sutiles. El esquema de normalización lineal por el máximo estructural de los sensores garantizó un aislamiento estricto entre particiones.

En concordancia con los criterios de transparencia y reproducibilidad de la Universidad de Cundinamarca, el código fuente en Python, el cuaderno interactivo de Google Colab con todas las salidas ejecutadas, los modelos entrenados (\texttt{modelo\_seleccionado.pt}) y los registros numéricos en CSV están disponibles públicamente en GitHub: \url{https://github.com/sofiadelcal/karen-redes-neuronales-iii-digitos}.

\clearpage\section{Referencias}
\refapa{Alpaydin, E. y Kaynak, C. (1998). Cascading classifiers. \textit{Kybernetika}, \textit{34}(4), 369--374.}
\refapa{Bishop, C. M. (2006). \textit{Pattern recognition and machine learning}. Springer.}
\refapa{Bottou, L. (2010). Large-scale machine learning with stochastic gradient descent. En Y. Lechevallier y G. Saporta (Eds.), \textit{Proceedings of COMPSTAT'2010} (pp. 177--186). Physica-Verlag. \url{https://doi.org/10.1007/978-3-7908-2604-3_16}}
\refapa{Glorot, X. y Bengio, Y. (2010). Understanding the difficulty of training deep feedforward neural networks. En \textit{Proceedings of the 13th International Conference on Artificial Intelligence and Statistics (AISTATS)} (Vol. 9, pp. 249--256). JMLR.}
\refapa{Goodfellow, I., Bengio, Y. y Courville, A. (2016). \textit{Deep learning}. MIT Press. \url{https://www.deeplearningbook.org/}}
\refapa{Grandini, M., Bagli, E. y Visani, G. (2020). Metrics for multi-class classification: An overview. \textit{arXiv preprint arXiv:2008.05756}. \url{https://doi.org/10.48550/arXiv.2008.05756}}
\refapa{Hastie, T., Tibshirani, R. y Friedman, J. (2009). \textit{The elements of statistical learning: Data mining, inference, and prediction} (2.ª ed.). Springer. \url{https://doi.org/10.1007/978-0-387-84858-7}}
\refapa{Kingma, D. P. y Ba, J. (2015). Adam: A method for stochastic optimization. \textit{3rd International Conference on Learning Representations (ICLR)}. \url{https://arxiv.org/abs/1412.6980}}
\refapa{Krogh, A. y Hertz, J. A. (1991). A simple weight decay can improve generalization. \textit{Advances in Neural Information Processing Systems}, \textit{4}, 950--957.}
\refapa{LeCun, Y., Bottou, L., Bengio, Y. y Haffner, P. (1998). Gradient-based learning applied to document recognition. \textit{Proceedings of the IEEE}, \textit{86}(11), 2278--2324. \url{https://doi.org/10.1109/5.726791}}
\refapa{Nair, V. y Hinton, G. E. (2010). Rectified linear units improve restricted boltzmann machines. En \textit{Proceedings of the 27th International Conference on Machine Learning (ICML)} (pp. 807--814).}
\refapa{Paszke, A., Gross, S., Massa, F., Lerer, A., Bradbury, J., Chanan, G., Killeen, T., Lin, Z., Gimelshtein, N., Antiga, L., Desmaison, A., Köpf, A., Yang, E., DeVito, Z., Raison, M., Tejani, A., Sasank, C., Vanschoren, J. y Chintala, S. (2019). PyTorch: An imperative style, high-performance deep learning library. \textit{Advances in Neural Information Processing Systems}, \textit{32}, 8024--8035.}
\refapa{Pedregosa, F., Varoquaux, G., Gramfort, A., Michel, V., Thirion, B., Grisel, O., Blondel, M., Prettenhofer, P., Weiss, R., Dubourg, V., Vanderplas, J., Passos, A., Cournapeau, D., Brucher, M., Perrot, M. y Duchesnay, É. (2011). Scikit-learn: Machine learning in Python. \textit{Journal of Machine Learning Research}, \textit{12}, 2825--2830.}
\refapa{Polyak, B. T. (1964). Some methods of speeding up the convergence of iteration methods. \textit{USSR Computational Mathematics and Mathematical Physics}, \textit{4}(5), 1--17. \url{https://doi.org/10.1016/0041-5553(64)90137-5}}
\refapa{Prechelt, L. (1998). Early stopping---but when? En G. Orr y K.-R. Müller (Eds.), \textit{Neural networks: Tricks of the trade} (pp. 55--69). Springer. \url{https://doi.org/10.1007/3-540-49430-8_3}}
\refapa{Ruder, S. (2016). An overview of gradient descent optimization algorithms. \textit{arXiv preprint arXiv:1609.04747}. \url{https://doi.org/10.48550/arXiv.1609.04747}}
\refapa{scikit-learn. (s.~f.). \textit{load\_digits: Optical recognition of handwritten digits dataset}. Scikit-learn documentation. \url{https://scikit-learn.org/stable/modules/generated/sklearn.datasets.load_digits.html}}
\refapa{Sokolova, M. y Lapalme, G. (2009). A systematic analysis of performance measures for classification tasks. \textit{Information Processing \& Management}, \textit{45}(4), 427--437. \url{https://doi.org/10.1016/j.ipm.2009.03.002}}
\refapa{Srivastava, N., Hinton, G., Krizhevsky, A., Sutskever, I. y Salakhutdinov, R. (2014). Dropout: A simple way to prevent neural networks from overfitting. \textit{Journal of Machine Learning Research}, \textit{15}(1), 1929--1958. \url{https://jmlr.org/papers/v15/srivastava14a.html}}
\refapa{Tieleman, T. y Hinton, G. (2012). Lecture 6.5-RMSprop: Divide the gradient by a running average of its recent magnitude. \textit{COURSERA: Neural Networks for Machine Learning}, \textit{4}(2), 26--31.}
\end{document}
